% ==========================================================================
% GENERATED by scripts/split_paper.py. Do not edit this file.
%
% Source: imported_papers/grokking/iclr2027_conference_v2.tex
% Part 12: Theory in tractable models
%
% Edits belong upstream, in the published paper. To change what the
% chapter includes, change imports/grokking.toml and re-run
% the splitter.
% ==========================================================================
\section{Theory in tractable models}
\label{app:theory}
This section works the mechanism through in models small enough to solve exactly, asking whether the measurements have to come out the way they do. The six results follow the order the main text needs them in, direction first, then how many units carry the circuit, then what the two ends cost:

\begin{itemize}\item A rule whose update carries a component orthogonal to the true gradient can still grok, as long as that transverse part dies out as the loss falls (Proposition~\ref{prop:persist}).
\item For the outgoing weight norms of a layer of $n$ units, $\mathrm{PR}$ (Equation~\ref{eq:pr}) runs from $1/n$, one unit doing everything, to $1$, every unit contributing equally, so $n \cdot \mathrm{PR}$ is the effective number of units carrying the circuit (Proposition~\ref{prop:pr}).
\item The credit knob changes how large each unit's step is, never which direction the network moves in (Proposition~\ref{prop:descent}).
\item Plain gradient flow keeps the ratios of unit sizes, so its solution inherits the initialization's participation. Giving every unit the same credit instead grows all units by one common amount, and participation rises monotonically with that growth, to $1$ in the small initialization regime the experiments run in (Proposition~\ref{prop:select}).
\item Give equal credit to \emph{groups} of units instead and the balancing happens only between groups. Inequality inside a group survives untouched, which is why only per-unit rules move the circuit (Proposition~\ref{prop:grain}).
\item Spreading the work over many units is what makes random deletion survivable. Deleting a fixed fraction of units removes, on average, the same fraction of every network's summed unit contributions. What differs is the variance around that average, a spread circuit has less of it, and the variance is what breaks the task (Proposition~\ref{prop:fluct}).
\end{itemize}

\begin{proposition}[Survival at the solution, not accumulated size, decides]
\label{prop:persist}
A rule ends training at a stationary point of its own update, which need not be a stationary point of the loss. Let $\vg = \nabla L$ and split the applied update $\vu_t$, the step the rule actually takes, into its part along the gradient and its part orthogonal to it,
\begin{equation}
\vu_t \;=\; \|\vg\|\,\Big(\kappa\,\frac{\vg}{\|\vg\|} \;+\; \sqrt{1-\kappa^{2}}\;\hat{\vb}_t\Big),
\qquad \hat{\vb}_t \perp \vg,\quad \|\hat{\vb}_t\| = 1.
\label{eq:mixedupdate}
\end{equation}
\begin{itemize}
\item[(i)] \emph{Fixed points.} Rescale the update to $\vg + \vb$ with $\vb \perp \vg$, so the parameters follow $\dot{\theta} = -\vg - \vb$. A stationary point satisfies $\nabla L = -\vb$, and orthogonality forces both sides to vanish separately. If $\vb(\theta^\star) = \vzero$ at every minimizer $\theta^\star$ of $L$, the stationary points of the applied dynamics are exactly those of $L$. If instead $\vb$ holds a fixed nonzero value $\vb_0$ there, take $L$ to include the weight decay term $\lambda\|\theta\|^{2}$, whose minimizers are isolated and whose Hessian $\mH = \nabla^{2}L(\theta^\star)$ is therefore invertible. Every minimizer is then displaced, to first order by $-\mH^{-1}\vb_0$. Along the directions the unregularized training loss leaves flat, the eigenvalues of $\mH$ are of order $\lambda$, so the displacement there is of order $\|\vb_0\|/\lambda$.
\item[(ii)] \emph{Descent.} Because $\hat{\vb}_t \perp \vg$ removes the cross term, $\langle \vu_t, \vg\rangle = \kappa\|\vg\|^{2} > 0$: for every $\kappa > 0$ and every choice of $\hat{\vb}_t$ the update is a descent direction, however large its orthogonal part.
\item[(iii)] \emph{Accumulation.} Write $\vb_t = \|\vg\|\sqrt{1-\kappa^2}\,\hat{\vb}_t$ for the transverse part, $\eta$ for the step size, and $\ve$ for a unit vector along a low curvature direction. The displacement accumulated along $\ve$ over $N$ steps, $\Delta = \eta \sum_{t=1}^{N} \langle \ve, -\vb_t\rangle$, grows as $N$ when $\hat{\vb}_t$ is fixed. It grows only as $\sqrt{N}$ when the $\hat{\vb}_t$ are independent and zero mean, since then $\mathbb{E}[\Delta]=0$ and $\mathrm{Var}[\Delta]=O(N)$.
\end{itemize}
\end{proposition}
\begin{proof}[Proof sketch]
(ii) and (iii) are the computations in the statement. For (i), a stationary point of $\dot{\theta} = -\vg - \vb$ satisfies $\nabla L = -\vb$, and $\langle \nabla L, \vb \rangle = 0$ makes both sides vanish separately: the stationary points are exactly the critical points of $L$ at which $\vb$ vanishes, which is all of them under the hypothesis. When instead $\vb$ holds the fixed value $\vb_0$ near $\theta^\star$, expanding $\nabla L(\theta^\star + \delta) = \mH\delta + O(\|\delta\|^2)$ and solving $\mH\delta = -\vb_0$ gives $\delta = -\mH^{-1}\vb_0$.
\end{proof}

Part (i) is the operative one: whether a rule groks is decided at the fixed points. The smallest eigenvalues of $\mH$ sit along the directions the training loss leaves flat, where only the weight decay curves the loss, so the displacement $\mH^{-1}\vb_0$ is large exactly there, the steep fitting directions barely move, and training accuracy survives when the circuit does not. Persistence is harmless whenever the component vanishes at the solution, since then $\vb_0 = \vzero$ however long the direction is held, which the error-gated control in Appendix~\ref{app:mech} confirms and which places a feedback controller, whose signal vanishes as its output reaches the target, on the tolerated side. Parts (ii) and (iii) say why size alone is not decisive: every mixture in Equation~\ref{eq:mixedupdate} descends, and only a fixed direction accumulates over the long plateau grokking traverses.

\paragraph{The two cases of part (i) cover different rules.} The case that displaces every minimizer is feedback alignment and the fixed bias control. A component that moves with the state falls under neither case, and the sign based optimizers are of that kind, so they are settled by convergence theory instead: raw sign updates provably fail to converge in general \citep{karimireddy2019errorfeedback} and sign-SGD memorizes here at every learning rate we swept, while Lion signs an interpolated momentum that provably converges \citep{chen2024lion} and Lion groks. The measurements agree with that split. Neither sign rule keeps a persistent transverse direction, at autocorrelation $-0.125$ for sign-SGD and $0.149$ for Lion over a lag of $4000$ epochs, while feedback alignment holds one at $0.937$, so the two memorizing routes separate: feedback alignment fails on fixed points and sign-SGD on convergence. Every rule in Table~\ref{tab:survival} is therefore scored against part (i), and the sign family by its convergence theory.
\begin{proposition}[Participation ratio bounds]
\label{prop:pr}
For $a \in \R^n_{\geq 0}$, $a \neq 0$, the participation ratio of Equation~\ref{eq:pr} satisfies $1/n \leq \mathrm{PR}(a) \leq 1$. The upper bound is attained exactly when all $a_j$ are equal, and the lower bound exactly when a single $a_j$ is nonzero. Consequently $n \cdot \mathrm{PR}(a)$ interpolates between $1$ and $n$ and acts as an effective count of contributing units.
\end{proposition}
\begin{proof}
By Cauchy--Schwarz, $(\sum_j a_j)^2 \leq n \sum_j a_j^2$ with equality exactly for constant $a$, which gives $\mathrm{PR} \leq 1$. Since the $a_j$ are nonnegative, $(\sum_j a_j)^2 \geq \sum_j a_j^2$ with equality exactly when at most one entry is nonzero, which gives $\mathrm{PR} \geq 1/n$.
\end{proof}

\begin{proposition}[The knob is a descent direction]
\label{prop:descent}
Let $g \neq 0$ be the gradient of the weight matrix the knob acts on, with one row $g_j$ per unit of the layer it feeds, and let $\tilde{g}$ be its rescaling by Equation~\ref{eq:knob} with any $\beta \in [0,1]$. Then $\langle \tilde{g}, g \rangle > 0$, so an update along $-\tilde{g}$ decreases the loss to first order.
\end{proposition}
\begin{proof}
Each row is rescaled by the positive scalar $s_j = \big((1-\beta)\|g_j\| + \beta \bar{g}\big) / \|g_j\|$, where $\bar{g} = \frac{1}{n}\sum_l \|g_l\|$ is the mean row norm of Equation~\ref{eq:knob} (rows with $\|g_j\| = 0$ are left at zero), so $\langle \tilde{g}, g \rangle = \sum_j s_j \|g_j\|^2 > 0$ whenever $g \neq 0$.
\end{proof}

\begin{proposition}[The two flows select different minima]
\label{prop:select}
Consider $n$ units in a diagonal linear network, where unit $i$ has one incoming weight $u_i$ and one outgoing weight $v_i$ and contributes $a_i = u_i v_i$, balanced so that $u_i = v_i > 0$ and $a_i = u_i^2$. Train by gradient flow from an initialization $a_i(0) = \epsilon_i > 0$ on the loss
\begin{equation}
L \;=\; \tfrac{1}{2}\Big(\sum\nolimits_i a_i - 1\Big)^{2},
\label{eq:toyloss}
\end{equation}
whose minima form the set $\{a : \sum_i a_i = 1\}$. The two flows converge to different members of that set.
\begin{itemize}
\item[(i)] Plain gradient flow keeps every ratio $a_i / a_j$ at its initial value, so it converges to the rescaled initialization $a_i = \epsilon_i / \sum_j \epsilon_j$, whose participation ratio equals that of the initialization.
\item[(ii)] The flow under the credit knob at $\beta = 1$ (Equation~\ref{eq:knob}) keeps every difference $u_i - u_j$ at its initial value, so it converges to $u_i = \sqrt{\epsilon_i} + C$ with $C$ set by $\sum_i (\sqrt{\epsilon_i} + C)^2 = 1$.
\item[(iii)] $\mathrm{PR}\big((\sqrt{\epsilon_i} + C)^2\big)$ is nondecreasing in $C \geq 0$, strictly increasing unless the $\epsilon_i$ are equal. Whenever the initialization undershoots the constraint, $\sum_i \epsilon_i < 1$ so that $C > 0$, the knob's solution therefore has strictly higher participation than the initialization, reaching $1$ as the initialization scale tends to zero. From an initialization that overshoots the constraint, $C < 0$ and participation can fall.
\end{itemize}
\end{proposition}
\begin{proof}
Write $r = \sum_j a_j - 1$ for the residual, a single scalar shared by every unit. The two flows differ in one line,
\begin{equation}
\text{plain flow: }\ \dot{u}_i = -\partial_{u_i} L = -r\,u_i,
\qquad\qquad
\text{knob: }\ \dot{u}_i = -r\,\bar{u},
\quad \bar{u} = \tfrac{1}{n}\textstyle\sum_j u_j,
\label{eq:twoflows}
\end{equation}
where the knob's form holds because unit $i$'s gradient row $(r u_i,\, r u_i)$ has norm $|r|\sqrt{2}\,u_i$, and Equation~\ref{eq:knob} at $\beta=1$ rescales it to the mean norm $|r|\sqrt{2}\,\bar{u}$. For (i), $\dot{a}_i = 2 u_i \dot{u}_i = -2 r a_i$, so $\frac{d}{dt} \log(a_i / a_j) = 0$ because $r$ is common to all units. The invariance is a property of this model, since in a network whose units see different residuals the ratios do move. For (ii), every $u_i$ takes the same step, so $\sqrt{a_i(t)} = \sqrt{\epsilon_i} + C(t)$ with $C(t) = -\int_0^t r \bar{u}\, ds$, and fitting the target fixes $C$. For (iii), write $b_i = \sqrt{\epsilon_i} + C$ and $S_k = \sum_i b_i^k$, so that $\mathrm{PR} = S_2^2 / (n S_4)$ and
\begin{equation}
\frac{d}{dC}\,\mathrm{PR} \;=\; \frac{4\,S_2\,\big(S_1 S_4 - S_2 S_3\big)}{n\,S_4^{2}},
\qquad
S_1 S_4 - S_2 S_3 \;=\; \tfrac{1}{2}\sum_{i,j} b_i b_j\,(b_i + b_j)(b_i - b_j)^2 \;\geq\; 0,
\label{eq:prmono}
\end{equation}
where the identity follows by symmetrizing the double sum $\sum_{i,j} b_i b_j^4 - b_i^2 b_j^3$ over $i$ and $j$, and it is zero exactly when all nonzero $b_i$ coincide. Participation is therefore nondecreasing along the path from $C = 0$, where it equals $\mathrm{PR}(\epsilon)$, to the fitted $C > 0$, and strictly increasing unless the $\epsilon_i$ are equal. As the initialization scale tends to zero, $C \to 1/\sqrt{n}$ and $a_i \to 1/n$, whose participation ratio is $1$. For $C < 0$ the bound does not apply, and simulation from overshooting initializations shows participation can fall.
\end{proof}

The contrast sits in Equation~\ref{eq:twoflows}. Plain flow is \emph{multiplicative}: each unit's step is proportional to its current size, a head start compounds, and the flow can only rescale the initialization's profile, never reshape it. The knob's flow is \emph{additive} in $\sqrt{a_i}$: every unit takes the same step, the head start survives only as a fixed offset that the common growth outgrows, and by part (iii) participation rises monotonically with that growth. Intermediate $\beta$ interpolates between the two. Simulating the flow at $n = 64$ over eight seeds confirms it: from initializations at participation $0.44$ and $0.085$, the fitted participation rises monotonically with $\beta$ and both reach the uniform solution at $\beta = 1$, with the loss driven to zero throughout, and the trained endpoints move by less than $0.01$ across initialization scales of $0.5$ to $4$ (Appendix~\ref{app:imposed}). This is the mechanism behind the knob sweep of Section~\ref{sec:select}: equalizing per-unit credit removes the dependence on the starting distribution.

Both flows minimize the same loss and arrive at different minima of it, which is the paper's claim in miniature: the task fixes the set of solutions, and the learning rule picks the member. Together with the small initialization $\ell_1$ bias of the product parameterization \citep{woodworth2020kernel}, the proposition says why plain flow preserves whatever concentration the early dynamics produce and why equalized credit overrides it. The model is linear where the trained networks are not, so we take the mechanism from the derivation and its confirmation from the trained ReLU networks of Section~\ref{sec:select}.

The same model derives the whole grain sweep, so the flat then steep curve of Section~\ref{sec:law} is a prediction rather than an observation. The result is Proposition~\ref{prop:select} run at two levels at once: between groups the flow is additive, exactly as the knob makes it between units, and within a group it stays multiplicative, exactly as plain flow.

\begin{proposition}[Group equalization removes only between group inequality]
\label{prop:grain}
Partition the $n$ units into $g$ groups of $m = n/g$. In the model of Proposition~\ref{prop:select}, rescale each unit's gradient row by the factor, common within its group, that equalizes total credit across groups. Write
\begin{equation}
p_i \;=\; \frac{\sqrt{\epsilon_i}}{\sum_{j \in \mathrm{group}(i)} \sqrt{\epsilon_j}}
\label{eq:groupshare}
\end{equation}
for unit $i$'s share of its own group's initialization, measured in the $\sqrt{a}$ coordinate the flow is linear in, and let $k$ run over groups. As the initialization scale tends to zero, the fitted solution has participation ratio
\begin{equation}
\mathrm{PR}(m) \;=\; \frac{S^2}{m\,Q},
\qquad
S = \frac{1}{g}\sum_{k}\sum_{i \in k} p_i^2,
\qquad
Q = \frac{1}{g}\sum_{k}\sum_{i \in k} p_i^4,
\label{eq:grainlaw}
\end{equation}
which equals the plain flow value of Proposition~\ref{prop:select} at $g=1$ and equals $1$ at $m=1$.
\end{proposition}

\begin{proof}
The rescaling factor is shared within a group, so inside group $k$ every $\log u_i$ has the same time derivative. The within group shares $p_i$ therefore freeze at their initial values, exactly as ratios freeze under plain flow. The group totals $U_k = \sum_{i \in k} u_i$ obey
\begin{equation}
\dot{U}_k \;=\; -\,r\,c_k U_k \;=\; -\,r\,\bar{U},
\qquad c_k = \bar{U}/U_k,
\qquad \bar{U} = \tfrac{1}{g}\textstyle\sum_k U_k,
\end{equation}
a common increment, so as the initialization tends to zero all group totals equalize to a common $U$. The solution is then $a_i = p_i^2 U^2$ with $U^2 = 1/(gS)$ forced by $\sum_i a_i = 1$, and substituting into Equation~\ref{eq:pr} gives Equation~\ref{eq:grainlaw}. At $m=1$ every share is $1$, so $S = Q = 1$ and $\mathrm{PR} = 1$, recovering Proposition~\ref{prop:select}. At $g=1$ the rescaling is the identity, and Equation~\ref{eq:grainlaw} reduces to the participation of the normalized initialization, the plain flow value.
\end{proof}

Coarse equalization therefore removes only the inequality \emph{between} groups. The inequality within a group is untouched, and while groups hold several units it dominates Equation~\ref{eq:grainlaw}, so participation stays near the plain flow value and can reach $1$ only when no group can hide any inequality, at one unit per group. That is the flat then steep shape. The comparison with the measured curve is parameter free: a lognormal initialization has plain flow participation $e^{-\sigma^2}$, so the measured $g=1$ anchor of $0.602$ fixes $\sigma$, and the law then predicts $0.674$, $0.737$, $0.834$ and $1.000$ at eight, four, two and one units per group, against the measured $0.623$, $0.695$, $0.753$ and $0.992$. It over predicts the interior, as expected when real credit magnitudes are correlated across units, and reproduces the flat regime, the location of the rise, and both anchors.

\begin{proposition}[Participation controls dropout fluctuation]
\label{prop:fluct}
Model the score of the correct class as a sum of unit contributions, $y = \sum_j c_j$ with $c_j \geq 0$ over $n$ units, and delete units independently with probability $q$, giving the surviving score $y_q = \sum_j b_j c_j$ with $b_j \sim \mathrm{Bernoulli}(1-q)$ marking survival. Then, for every failure margin $\varepsilon \in (0,1)$,
\begin{equation}
\frac{\sqrt{\mathrm{Var}[y_q]}}{\mathbb{E}[y_q]}
= \sqrt{\frac{q}{1-q}} \cdot \frac{1}{\sqrt{n \cdot \mathrm{PR}(c)}},
\qquad
\mathbb{P}\big(y_q < (1-\varepsilon)\,\mathbb{E}[y_q]\big)
\;\leq\; \frac{q}{(1-q)\,\varepsilon^2 \, n \cdot \mathrm{PR}(c)},
\label{eq:fluct}
\end{equation}
where $\mathrm{PR}(c)$ is Equation~\ref{eq:pr} evaluated on the vector of contributions.
\end{proposition}
\begin{proof}
$\mathbb{E}[y_q] = (1-q) \sum_j c_j$ and, by independence, $\mathrm{Var}[y_q] = q(1-q) \sum_j c_j^2$. Dividing and using Equation~\ref{eq:pr} gives the first identity, and Chebyshev's inequality gives the tail bound.
\end{proof}

Deletion removes, on average, the same fraction of every circuit's summed contributions, so what separates circuits is the fluctuation around that mean, and Proposition~\ref{prop:fluct} says the fluctuation falls with the effective unit count $n \cdot \mathrm{PR}$. The contributions $c_j$ are not directly observable, and circuit participation, the $\mathrm{PR}$ of the outgoing norms, is the measured stand-in. With backpropagation's circuit at about $154$ effective units and the knob's at $254$, the failure bound for the concentrated circuit is about $1.6$ times the distributed one at every deletion rate, the ordering measured in Tables~\ref{tab:targeted}, \ref{tab:dissoc} and~\ref{tab:parity}. The model simplifies in three ways: the network combines unit outputs through a second layer and a softmax rather than a plain sum, contributions to competing classes can cancel, and the accuracy a task needs is not a fixed fraction of the mean score. We therefore take from it the choice of measure and the direction of the effect, and take magnitudes from the experiments.
